% Junior Level 3, Week 9 -- answer key
\answersweek{9}{Project Studio}{Answers for workbook chapter 9 and play-at-home card 9.}

\section{Lesson questions}
\begin{answer}[Exit question]
The child points at the code for that part. Scratch: the script or the blocks that make it, for example the repeat inside a repeat that draws the flower, or \enquote{change flowers by 1} for the counter. Python: the lines that make it, for example the \texttt{for} loop that draws the roof and the indented lines under it. Pointing at the right place and saying one sentence about it counts.
\end{answer}

\begin{answer}[Reflection]
Any number of colored parts is fine. Three parts and a finished partner check means milestone M2 is reached. Fewer parts: plan extra studio time (buffer B2 or the first 15 minutes of week 10). Any honest smiley is right.
\end{answer}

\begin{answer}[Can everyone see it? (safety talk)]
Any of the four tips is a good answer: thick lines (pen size 5 or more), dark next to light (not yellow on white), not color alone (add a word or a shape), and big enough sprites, shapes and words. A project passes the squint test when every part can still be seen with squinted eyes from a step back.
\end{answer}

\section{Workbook exercises}
\begin{answer}[Exercise 9.1: My Plan Card (milestone M1)]
Every plan card is different. Look for: one path circled; a drawing of the finished project; three parts with one idea each, in building order (Scratch example: get ready, the flower pattern, the counter; Python example: the house, the roof, a sun); the must-haves checked (Path A: a repeat inside a repeat or a My Block, plus a score or a counter; Path B: a \texttt{for} loop, at least two colors and one shape changed or added); and a credit line when an idea was borrowed (\enquote{the starter}, \enquote{Kai's flower}). Too many ideas: move the extras to the back of the card.
\end{answer}

\begin{answer}[Exercise 9.2: Build My Big Build (milestone M2)]
\textbf{Leo's Flower Garden (Path A):} the green flag clears the stage and sets the flowers counter to 0. Each press of the space key draws a flower of six squares where the mouse-pointer is, and the counter on the stage goes up by 1.

\textbf{Omar's postcard (Path B):} the starter's house is now blue and bigger (sides of 120 instead of 100), with a gold roof. The roof still sits on the house because all three \texttt{100}s were changed together.

\textbf{What to look for in every project:} \tickbox{} each part tested before the next one was started \quad \tickbox{} the must-haves for the chosen path \quad \tickbox{} it starts with the green flag or Run \quad \tickbox{} it was saved. Full reference: \texttt{exercises/week09/ex-my-big-build/solution.md}, the project teacher guide and the reference postcard.
\end{answer}

\begin{answer}[Exercise 9.3: Partner Check]
All six rows say \enquote{yes} for a finished project. A \enquote{not yet} is a good find, not a failure: the owner fixes it next. Look for a bug report in the \enquote{I expected \ldots, but it \ldots} shape, for example \enquote{I expected the counter to go back to 0, but it kept the old number.} Common fixes: a hat block on top of a script that never starts; \enquote{set \ldots{} to 0} under the green flag; a Python typo (\texttt{t.colr}, \texttt{"bleu"}, a missing \texttt{)} or quote mark); an indent for a line that belongs inside the \texttt{for} loop; thicker lines or darker colors so everyone can see it.
\end{answer}

\section{Play-at-home card}
\begin{answer}[Practice the show]
Look for three short sentences in the order of the plan card, one part explained in the child's own words (\enquote{This loop draws the roof. It goes three times, because a triangle has three sides.}) and one change they made (\enquote{I made the house blue}).
\end{answer}

\begin{answer}[The squint test]
Any drawing is right. Good ideas to try: thicker lines, darker colors next to light ones, bigger shapes, or a word as well as a color.
\end{answer}
